\begin{figure}[tbp]
\centering
\begin{tikzpicture}[
  box/.style={draw,rounded corners=2pt,align=center,text width=4.7cm,
    inner sep=6pt,font=\small},
  arr/.style={-{Stealth[length=2mm]},thick},
  node distance=10mm and 13mm]
\node[box,text width=8.2cm] (engine)
  {Simultaneous pruning in actual simple groups\\
   and the protected-path bound};
\node[box] (affine) at (-3.2,-1.7)
  {Private affine charts:\\a nonrepresentable family};
\node[box] (marker) at (3.2,-1.7)
  {Reverse tests and\\the forward marker theorem};
\node[box,below=of affine] (ceiling)
  {Effective parameter control:\\a computable size ceiling\\for one negative lattice};
\node[box,below=of marker] (weighted)
  {Weighted representations:\\one cell alignment for\\all arithmetic readings};
\node[box,below=of ceiling] (graph)
  {Finite graph witnesses:\\affirmative search and\\conditional size bounds};
\node[box,below=of weighted] (count)
  {Independent template count:\\lattice size bounded by units,\\not numerical weights};
\node[box,below=of graph] (minimum)
  {An ordinary Busy Beaver cutoff\\and finite selection of a\\minimum negative order table};
\node[box,below=of count] (undec)
  {Bounded replacement of\\Diophantine solutions:\\no representation decider};
\draw[arr] (engine.south west) -- (affine.north);
\draw[arr] (engine.south east) -- (marker.north);
\draw[arr] (affine) -- (ceiling);
\draw[arr] (marker) -- (weighted);
\draw[arr] (ceiling.west) -- ++(-0.35,0) |- (minimum.west);
\draw[arr] (graph) -- (minimum);
\draw[arr] (graph.south east) -- (undec.north west);
\draw[arr] (count) -- (undec);
\draw[arr] (weighted.east) -- ++(0.35,0) |- (undec.east);
\end{tikzpicture}
\caption{The two uses of the structural argument.  The affine branch
gives a negative lattice and a computable size ceiling independently of
any hypothetical decider.  The weighted branch combines simultaneous
arithmetic readings with a separate, weight-independent count; a
decider would then bound positive Diophantine solutions.  Finite graph
witnesses also supply the bounded searches used, with the size ceiling,
in the minimum-table specification.}
\label{fig:proof-overview}
\end{figure}
